\section{密度算符与一体密度矩阵}

Gross 第 6 章把“在 $\mathbf{r}$ 找到粒子的概率”提升为二次量子化算符，并引入一体密度矩阵。这是等离子体、HF 与 DFT 的共同语言。

\subsection{一体密度算符}

$N$ 粒子体系中，
\begin{equation}
  \hat n(\mathbf{r})
  =\sum_{i=1}^N \delta(\mathbf{r}-\hat{\mathbf{r}}_i)
  =\sum_\sigma\psi_\sigma^\dagger(\mathbf{r})\psi_\sigma(\mathbf{r}).
\end{equation}
期望值 $n(\mathbf{r})=\langle\hat n(\mathbf{r})\rangle$。Fourier 分量
\begin{equation}
  \hat n_{\mathbf{q}}
  =\int\mathrm{d}\mathbf{r}\,e^{-\mathrm{i}\mathbf{q}\cdot\mathbf{r}}\hat n(\mathbf{r})
  =\sum_{\mathbf{k}\sigma}
  c_{\mathbf{k}\sigma}^\dagger c_{\mathbf{k}+\mathbf{q},\sigma}
\end{equation}
直接进入库仑相互作用 $\sum_{\mathbf{q}}v_{\mathbf{q}}(\hat n_{\mathbf{q}}\hat n_{-\mathbf{q}}-N)/2L^d$。

\subsection{一体密度矩阵}

\begin{equation}
  \rho(\mathbf{r}\sigma,\mathbf{r}'\sigma')
  =\bigl\langle
  \psi_{\sigma'}^\dagger(\mathbf{r}')\psi_\sigma(\mathbf{r})
  \bigr\rangle.
\end{equation}
对角元即密度；非对角元编码相干与交换。对单一 Slater 行列式，
\begin{equation}
  \rho(\mathbf{x},\mathbf{x}')
  =\sum_i^{\mathrm{occ}}
  \phi_i(\mathbf{x})\phi_i^*(\mathbf{x}'),
  \qquad
  \rho^2=\rho
\end{equation}
（幂等性是“纯行列式”的充要条件，Gross/L\"owdin）。自旋分辨的 $\rho_{\uparrow},\rho_{\downarrow}$ 决定 HF 交换核。

\subsection{与对关联}

二体密度
\begin{equation}
  n_2(\mathbf{r},\mathbf{r}')
  =\bigl\langle
  \psi^\dagger(\mathbf{r})\psi^\dagger(\mathbf{r}')\psi(\mathbf{r}')\psi(\mathbf{r})
  \bigr\rangle
  =n(\mathbf{r})n(\mathbf{r}')g(\mathbf{r},\mathbf{r}'),
\end{equation}
$g$ 即前文对关联函数。势能完全由 $n_2$ 或等价地由 $S(\mathbf{q})$ 决定。
